% Fuente: Auxiliar 3, sección «Poliedro».
\begin{enumerate}
    \item[(a)] \textbf{Verdadera.}

    Si \(n=m+1\), entonces el espacio afín
    \[
    \{x\in\mathbb{R}^n\mid Ax=b\}
    \]
    tiene dimensión
    \[
    n-m=1.
    \]
    Esto quiere decir que es una recta, si consideramos la condición \(x\geq 0\), \(P\) puede ser vacía, un punto, un segmento o un rayo.

    Las soluciones básicas factibles corresponden a puntos extremos de \(P\).
    Como un subconjunto convexo de una recta puede tener a lo más dos puntos extremos, entonces \(P\) tiene a lo más dos soluciones básicas factibles distintas.

    \item[(b)] \textbf{Falsa.}

    El conjunto de soluciones óptimas no necesariamente es acotado.


    Contraejemplo: Considere el siguieente LP.
    \[
    \min \, x_1 - x_2 \quad \text{sujeto a } x_1 - x_2 = 0,\quad x_1,x_2\ge 0.
    \]
    La restricción \(x_1 - x_2 = 0\) fuerza a que la función objetivo valga \(0\) en todo punto factible, por lo que el conjunto de soluciones óptimas coincide con la región factible \(\{(t,t) : t\ge 0\}\), un rayo no acotado.
    
    Es posible construir el contrajemplo de forma matricial (otra forma de responder)
        \[
    P=\{x\in\mathbb{R}^2\mid x_1-x_2=0,\; x_1,x_2\geq 0\}.
    \]
    Por lo tanto,
    \[
    P=\{(t,t)\mid t\geq 0\},
    \]
    considerando la función objetivo $\min c^Tx$ con $c=(1,-1)$, todos los puntos de P son óptimos, y como P es un rayo no acotado, existen infinitas soluciones.

    \item[(c)] \textbf{Falsa.}

    Es cierto que una solución básica factible tiene a lo más \(m\) variables básicas, y por lo tanto a lo más \(m\) variables positivas. Sin embargo, una solución óptima no necesariamente debe ser básica.

    Considerando el mismo contraejemplo de (b), $m=1$, cualquier punto del conjunto $\{(t,t)\:|\:t>0\}$ pertenece a P, es óptimo y tiene 2 variables positivas.


    \item[(d)] \textbf{Verdadera.}

    Suponga que existen dos soluciones óptimas distintas \(x^1\) y \(x^2\), con
    \[
    x^1\neq x^2.
    \]
    Como \(P\) es convexo, para todo \(\lambda\in[0,1]\),
    \[
    x^\lambda=\lambda x^1+(1-\lambda)x^2\in P.
    \]

    Además, como la función objetivo de un problema lineal es lineal, se tiene
    \[
    c^\top x^\lambda
    =
    c^\top\left(\lambda x^1+(1-\lambda)x^2\right).
    \]
    Luego,
    \[
    c^\top x^\lambda
    =
    \lambda c^\top x^1+(1-\lambda)c^\top x^2.
    \]
    Como \(x^1\) y \(x^2\) son óptimos,
    \[
    c^\top x^1=c^\top x^2=z^\star.
    \]
    Por lo tanto,
    \[
    c^\top x^\lambda
    =
    \lambda z^\star+(1-\lambda)z^\star
    =
    z^\star.
    \]
    Así, todo punto del segmento
    \[
    \{\lambda x^1+(1-\lambda)x^2\mid \lambda\in[0,1]\}
    \]
    es óptimo.

    Como el intervalo \([0,1]\) contiene incontablemente muchos puntos, existen incontablemente muchas soluciones óptimas.

    \item[(e)] \textbf{Falsa.}

    No necesariamente existen dos soluciones básicas factibles óptimas.

    Considere nuevamente el contraejemplo en (b), el conjunto óptimo es $\{(t,t)\mid t\geq 0\}$, hay infinitas soluciones óptimas. Sin embargo, el conjunto \(P\) es un rayo cuyo único punto extremo es $(0,0)$
    Por lo tanto, existe una única solución básica factible óptima.

    Así, puede haber múltiples soluciones óptimas sin que existan al menos dos soluciones básicas factibles óptimas.
    
\end{enumerate}
